% SPDX-License-Identifier: Apache-2.0
\section{Implizite Erledigung und direkter Prüfungszugang}
\label{sec:assessment}
\begin{principle}{Verbindliche Produktentscheidung vom 31. Oktober 2026}
Jedes gemeisterte Lernziel beendet seinen eigenen Vorbereitungsbedarf.
Ein von keinem offenen Ziel mehr benötigtes Nicht-Prüfungsziel ist implizit
erledigt und beendet ebenfalls seinen Bedarf. Jede Abschlussprüfung bleibt
ohne Mastery offen und muss direkt ansteuerbar sein. Manuelle
Freistellungsprofile und zusätzliche Prüfbelege sind dafür nicht erforderlich.
\end{principle}

Mastery bleibt der einzige gespeicherte Ergebniszustand pro atomarem Lernziel.
Implizite Erledigung wird abgeleitet und niemals als Mastery zurückgeschrieben.
Grundlage ist die \srcref{docs/concept/skill-graph/assessment-gated-completion.md}{aktuelle
Konzeptspezifikation}; der \srcref{docs/dev/assessment-gated-completion-pilot.md}{Implementierungsstand}
trennt lokale Tests, Deployment und tatsächliche Coach-Abnahme.

\subsection{Fester Umfang und direkte Bedarfsrelation}
Für eine Curriculumgeneration und das verbindliche Personal Curriculum sei
$T$ die Menge der eindeutigen atomaren Ziele des gewählten Umfangs.
Maßgeblich ist die aufgelöste Program-/Kompositionsstruktur. Ein temporärer
Lernfokus oder eingeklappter Baumzweig verändert diesen Umfang nicht.
Die Standardauswertung umfasst das Personal Curriculum; ein ausdrücklich
gewählter Abschnitt erhält seine eigene vollständige Auswertung.

Cluster-Voraussetzungen werden in atomare Pflichten expandiert;
vererbte Anforderungen und autoritative atomare \code{goalEntry}-Projektionen
werden berücksichtigt. Die vollständige Voraussetzungshülle $V_H$ enthält
auch versteckte und fachübergreifende Unterstützung. Unterstützende Ziele
werden dadurch nicht zu Zielen des Abschlussnenners.

Die Relation
\begin{equation}
 F_H=\{(u,v)\mid u\text{ benötigt unmittelbar }v
       \text{ nach atomarer Normalisierung}\}
\end{equation}
zeigt vom Verbraucher zur Voraussetzung. Fehlende Referenzen, Identitätsfehler
sowie Zyklen in \code{contains}, \code{requires} oder der expandierten
atomaren Relation sind Berechnungsfehler. Sie bedeuten nicht, dass keine
Arbeit mehr erforderlich sei.

\textbf{Eine statische transitive Hülle ist hier falsch.}
Bei $C\to B\to A$ beendet Mastery von B den Bedarf an A.
Eine zusätzliche direkte Kante $C\to A$ wäre dagegen eine eigenständige
Anforderung von C. Transitive Hülle und transitive Reduktion sind daher
keine bedeutungserhaltenden Umformungen dieser Abschlusssemantik.

Sei $Q_H\subseteq V_H$ die autoritativ bestimmte Menge der Abschlussprüfungen.
Die aktuelle Abbildung nutzt die bestehenden Prüfungs-/Assessmenttypen,
nicht Titel oder \code{coveredGoalIds}. Bereits im Kontext enthaltene
Prüfungspflichten verschwinden nicht dadurch, dass sie nur als Unterstützung
oder in einer anderen Ansicht sichtbar sind. Alternative Aufgabenformen
und Wahlentscheidungen müssen bereits bei der Umfangsauflösung geklärt sein.

\subsection{Mastery, Bedarf und kleinster Fixpunkt}
Mit der aktuellen globalen Mastery-Menge $M$ sei $E=M\cap V_H$.
Für eine Menge $X$ bereits ausdrücklich oder implizit erledigter Ziele gilt
\begin{align}
 \Consumers_H(v)&=\{u\mid(u,v)\in F_H\},\\
 \Active_H(X)&=\{(u,v)\in F_H\mid u\notin X\}.
\end{align}
Mastery beendet die \emph{eigenen ausgehenden} Anforderungen. Ein eingehender
Bedarf an einem bereits gemeisterten Ziel wird nicht durch dieses Ziel
hindurch an dessen frühere Vorbereitung weitergereicht.

Die drei Produktregeln ergeben unmittelbar
\begin{align}
 X_0&=E,\\
 X_{n+1}&=X_n\cup\{v\in V_H\setminus Q_H\mid
              \Consumers_H(v)\subseteq X_n\},\label{eq:waiver-step}\\
 X_*&=\bigcup_{n\geq0}X_n,\qquad I_H=X_*\setminus E.
 \label{eq:waiver-result}
\end{align}
Der Operator ist monoton und inflationär. Auf dem endlichen Graphen erreicht
die Iteration ihren kleinsten Fixpunkt. Dabei gilt stets
$I_H\cap E=\varnothing$ und $I_H\cap Q_H=\varnothing$.
Es entsteht keine zusätzliche gespeicherte Ergebnisart.

Die leere Verbrauchermenge ist kein Sonderfall: Ein gewöhnliches Endziel
ohne Verbraucher ist implizit erledigt. Dies gilt auch für Orientierung
und Gedächtnisziele. Die ausdrücklichen Interaktions- und Recall-Verträge
bleiben unverändert; weder ein Kartenerfolg noch eine motivational erfolgte
Interaktion wird durch die Auswertung erfunden. Enthält ein gültiger endlicher
Umfang keine Abschlussprüfung, wird seine gesamte nicht gemeisterte Arbeit
implizit erledigt. Zusätzliche Ausnahmen für einzelne Zielarten sind in
diesen drei Regeln nicht vorgesehen.

\paragraph{Äquivalenter linearer Algorithmus.}
Starte bei allen noch nicht gemeisterten Abschlussprüfungen. Folge ihren direkten
Voraussetzungskanten, besuche jedes Ziel höchstens einmal und stoppe an
\emph{jedem} gemeisterten Ziel. Die erreichten nicht gemeisterten Ziele
bleiben benötigt. Alle übrigen nicht gemeisterten Ziele sind implizit erledigt.
Auf einem DAG entspricht dies genau dem Fixpunkt und benötigt
$O(|V_H|+|F_H|)$ Zeit und Speicher. Eine quadratische Speicherung aller
indirekten Beziehungen entfällt.

\subsection{Beispiel: Mastery unterbricht den Bedarf}
\begin{figure}[htbp]
\centering
\begin{tikzpicture}[node distance=12mm]
\node[goal,fill=orange!12] (c) at (0,0) {C: Prüfung\\offen};
\node[goal,fill=green!9] (b) at (4.5,0) {B: Mastery};
\node[goal] (a) at (9,0) {A: implizit};
\draw[req,draw=warm,dashed] (c)--node[above,info]{benötigt B}(b);
\draw[req] (b)--node[above,info]{Bedarf endet}(a);
\node[info,text width=115mm] at (4.5,-1.4)
 {C bleibt offen. B ist schon gemeistert und benötigt A nicht mehr.\\
 Ohne weiteren offenen Verbraucher ist A implizit erledigt.};
\end{tikzpicture}
\caption{Mastery beendet Vorbereitung unabhängig von einem Prüfungserfolg.
Eine zusätzliche direkte Kante von C nach A würde A weiterhin erforderlich machen.}
\label{fig:waivers}
\end{figure}
\FloatBarrier

\begin{table}[htbp]
\centering
\begin{tabularx}{\linewidth}{@{}lX@{}}
\toprule
Fall & Ergebnis\\
\midrule
$C\to B\to A$, nur B gemeistert & A implizit; Prüfung C bleibt offen.\\
Zusätzlich $C\to A$ & A bleibt direkt für C erforderlich.\\
Prüfungen B und C benötigen A & Ein einzelner Erfolg reicht nicht, solange der andere Verbraucher offen ist.\\
Beide Prüfungen gemeistert & A und anschließend seine nicht mehr benötigte Vorbereitung werden implizit erledigt.\\
Gewöhnliches Endziel ohne Verbraucher & Implizit erledigt, auch ohne bisherigen Erfolg.\\
Prüfung ohne Verbraucher & Bleibt bis zu ihrer Mastery offen.\\
\bottomrule
\end{tabularx}
\caption{Bedarfsrechnung ohne manuelle Freistellungsregeln.}
\label{tab:waivers}
\end{table}

\subsection{Abschluss und direkte Prüfungsauswahl}
Der Umfang ist abgeschlossen genau dann, wenn
\begin{equation}
 \Complete(H)\iff T\subseteq X_*\ \land\ Q_H\subseteq E.
\end{equation}
Auf dem endlichen DAG sind alle Prüfungspflichten gemeistert auch hinreichend:
Ohne offene Prüfung bleibt keine dauerhafte Nachfrage nach Nicht-Prüfungsarbeit.
Der Nenner zählt trotzdem nur Ziele aus $T$, nicht zusätzliche Unterstützung.
Eine in Sekundarstufe~I implizit erledigte Kompetenz kann im größeren
stufenübergreifenden Umfang noch benötigt werden; ihre globale Mastery bleibt gleich.

Jede passende Abschlussprüfung ist ohne vorherige Voraussetzungsmastery direkt
wählbar. Bewusste Auswahl, Authentifizierung, Lernendenrevision,
Sitzungszuordnung und verfügbare Aufgaben bleiben erforderlich. Ein Fehlversuch
erzeugt keinen negativen Ergebnisdatensatz und nimmt einen früheren Erfolg
nicht zurück. Ausführliche KI-Bewertung mit Punkten, Begründungen und
Verbesserungshinweisen bleibt im Gespräch erhalten.

Die lesende Übersicht schreibt keine Navigation. Suche und seitenweise
Darstellung erschließen das gesamte Angebot; eine Seitengrenze darf keine
Prüfung dauerhaft unzugänglich machen. Empfehlungen dürfen den aktuellen
Abschnitt bevorzugen, ohne die verbindlichen Pflichten zu verkleinern.
Normale Pflichtvorschläge lassen implizit erledigte Arbeit aus. Freiwilliges
Lernen bleibt über die unveränderte fachliche Zugangsprüfung möglich.

\subsection{Cockpit, Lernplan und technische Grenzen}
Der Baum verwendet ein kompaktes Icon mit dem Tooltip
„Für diesen Abschnitt nicht mehr erforderlich“. Diese Aussage behauptet
keinen Prüfungserfolg, wenn ein gewöhnliches Lernziel den Bedarf beendet hat.
Mastery-Anzeigen und abgeleitete Erledigung bleiben unterscheidbar.

Für die feste, eindeutige Planbasis $G$ gilt standardmäßig
\begin{equation}
 I_{\mathrm{plan}}=|G\cap(M\cup I_H)|.
\end{equation}
Die automatische Bedarfsrechnung verwendet das vollständige Personal Curriculum;
erst anschließend wird mit der Planmenge $G$ geschnitten. Ein datierter Plan
ist keine eigene abgeschlossene Prüfungswelt: Ein reiner Vorbereitungsplan
darf nicht schon deshalb fertig sein, weil seine Prüfung außerhalb von $G$
liegt. Auch echte fachübergreifende Anforderungen bleiben Bedarf, während
Quoten und Fachbilanzen getrennt bleiben. Eine kleine Abschnittsauswertung
ändert die Plananrechnung nicht.
Die Cockpit-Einstellung „Nur explizit erreichte Lernziele zählen“ ist
standardmäßig aus; bei Auswahl zählt nur $|G\cap M|$.
Planbasis, Termine und Sollmengen ändern sich dadurch nicht. Der Periodenwert
$H$ zählt weiterhin ausschließlich tatsächliche, zeitlich belegte Abschlüsse;
aktuelle implizite Erledigung erzeugt keine historischen Lernerfolgsereignisse.

„Prüfungsabkürzungen anbieten“ steuert direkte Einstiege. Die Einstellung
verändert weder Mastery noch den berechneten Bedarf oder gültige Plananrechnung.
Der \srcref{docs/concept/didactic/unified-learning-plan-status.md}{gemeinsame
Lernplanstatusvertrag} gilt für Cockpit und Coach gleichermaßen.

Statische Kompilierungen werden begrenzt nach Curriculumgeneration,
Projektion, Zielumfang und Semantikrevision zwischengespeichert. Mastery und
implizite Ergebnisse sind keine zweite Wahrheitsquelle. Nach Ergebnis- oder
Umfangsänderungen wird neu gerechnet; alte implizite Mengen werden niemals
zu Mastery-Eingaben. Ungültige Daten bleiben als Einschränkung sichtbar.
Automatische Tests prüfen diese Verträge, ersetzen aber weder Aufgabenqualität
noch menschliche Erprobung oder nachgewiesene Lernwirksamkeit.
